\section{Whole minimum-degree vertices at order \texorpdfstring{$2\delta+3$}{2delta+3}}
\label{sec:wholecore}

Throughout this section, $D$ is strongly connected, arc-minimal and
all-deficient,
$n=2\delta+3$, and $\delta=\delta^+(D)$. Let $z$ count whole
minimum-outdegree vertices, and let $h$ count vertices with $g(x)=2$.
We will prove
\begin{equation}\label{eq:wholecore}
z\ge h+7.
\end{equation}


\subsection{Degree balance at the first odd boundary}
\begin{corollary}\label{cor:near}
At $n=2\delta+3$, let $r_i$ count vertices with outdegree $\delta+i$.
Then $i\in\{0,1,2\}$ and
\[
 r_0=M+r_2\ge\delta+5+r_2.
\]
Thus a strict majority of the vertices have minimum outdegree.
\end{corollary}
\begin{proof}
The gap bound gives $2d^+(x)-g(x)\le n-1$ and $g(x)\le2$, so all
outdegrees are at most $\delta+2$. Counting arcs gives $M=r_0-r_2$.
Use the uniform margin in Theorem~\ref{thm:residual-five-excluded}.
\end{proof}

For each target $x$, use $T_x,P_x$ from the foundational target-packet section and write
\[
e_x=\dout(x)-\delta,\quad q_x=|T_x|-\delta=e_x+\mu(x),\quad
p_x=|P_x|.
\]
Counting outgoing arcs from $P_x$ gives, whenever $p_x>0$,
\begin{equation}\label{eq:capacity}
\delta p_x+\sum_{v\in P_x}e_v
\le \binom{p_x}{2}+p_x(\delta+q_x-p_x).
\end{equation}
In particular, $q_x=0$ implies $p_x=0$, and $q_x>0$ implies
$p_x\le2q_x-1$. Define
\[
\sigma_x=
\begin{cases}
2q_x-1-p_x,&q_x>0,\\
0,&q_x=0.
\end{cases}
\]
Let $K=\{x:\sigma_x>0\}$, $k=|K|$, and
$\Sigma=\sum_x\sigma_x$. A target in $K$ with $\sigma_x=1$ is called
a \emph{unit-slack target}; a target with $\sigma_x\ge2$ is called
a \emph{higher-slack target}.

Put $U(v)=V(D)\setminus(\{v\}\cup\Nout(v)\cup\Ntwo(v))$.
Lemma~\ref{lem:gap} gives
\begin{equation}\label{eq:excess}
|U(v)|=2-2e_v+g(v),\qquad e_v\in\{0,1,2\},
\end{equation}
and $e_v=2$ forces $g(v)=2$ and $U(v)=\varnothing$.
If $r_i$ counts excess-$i$ vertices, then $M=r_0-r_2$.
Counting far source-target pairs and total degrees gives
\[
\sum_xp_x=2M+n+h,\qquad \sum_xq_x=n+M.
\]
Since the $q_x=0$ targets are precisely the $z$ whole minimum-degree
vertices,
\begin{equation}\label{eq:slacksum}
{\Sigma=z-h.}
\end{equation}

\begin{lemma}\label{lem:zeroslack}
A zero-slack target has exactly one of the following forms:
\[
q_x=0,\quad e_x=\mu(x)=p_x=0;
\qquad\text{or}\qquad
q_x=1,\quad e_x=0,\quad \mu(x)=p_x=1.
\]
In the second case, the unique vertex of $P_x$ has minimum outdegree.
Every source is far from at most one zero-slack target.
\end{lemma}
\begin{proof}
Suppose $q_x>0$ and $\sigma_x=0$. Equality holds throughout
\eqref{eq:capacity}: all vertices of $P_x$ have outdegree $\delta$,
$D[P_x]$ is a regular tournament, and $P_x$ completely dominates
$R_x=T_x\setminus P_x$.
If $q_x\ge2$, then $|P_x|=2q_x-1\ge3$ and $P_x$ has an internal arc,
contradicting Lemma~\ref{lem:uniform}.
Hence $q_x=p_x=1$. The alternative $e_x=1,\mu(x)=0$ contradicts
Proposition~\ref{prop:whole}, so $e_x=0,\mu(x)=1$.

If $P_x=\{v\}$, capacity equality gives
$T_x=\{v\}\dot\cup\Nout(v)$.
For any other target $y\in U(v)$, we have $y\notin T_x$, so $y\to x$.
Two zero-slack far targets of the same source would thus form a digon.
\end{proof}

It follows that every nonminimum-degree vertex belongs to $K$.
Every minimum-degree vertex has three or four far targets by
\eqref{eq:excess}; at most one is outside $K$, so at least two lie in $K$.
In particular, $K$ is nonempty.

\subsection{Protected predecessors and higher-slack targets}

\begin{lemma}\label{lem:parent}
Every unit-slack target $x$ has $e_x\le1$. If $p_x>0$, it has a
vertex $y\in K$ such that
\[
y\to x,\qquad q_y>q_x,\qquad P_x\cup\{x\}\subseteq P_y.
\]
\end{lemma}
\begin{proof}
For $q_x>0$, the residual identity reads
$\eta(x)=\sigma_x+g(x)-2e_x$.
Thus $\sigma_x=1$ and $\eta(x)\ge0$ imply $2e_x\le3$.

If $p_x>0$, then $q_x\ge2$, $p_x=2q_x-2$, and
$|R_x|=\delta-q_x+2$. Omitting even one vertex of $R_x$ from
$\bd P_x$ would bound the total outdegree of $P_x$ by
$(\delta-\tfrac12)p_x<\delta p_x$. Hence $\bd P_x=R_x$.
Now $|W_x|=|R_x|$ for $W_x=\Nin(x)$, and
$X_x=\bd^2P_x\subseteq W_x$ has size less than $|R_x|$.
Choose $y\in W_x\setminus X_x$. No arc from $T_x=P_x\cup R_x$
points to $y$, so
\[
\Nin(y)\subseteq\Nin(x)\setminus\{y\}.
\]
Consequently $q_y>q_x$, and $x$ cannot reach $y$ in at most two steps.
Nor can any vertex of $P_x$ reach $y$: an arc to $y$ would give a
two-step path to $x$, and a two-step path to $y$ would pass through
$\Nin(y)\subseteq\Nin(x)$.
This proves the inclusion. Since $q_y>q_x\ge2$,
Lemma~\ref{lem:zeroslack} also gives $y\in K$.
\end{proof}

\begin{lemma}\label{lem:twoheavy}
There are at least two higher-slack targets.
\end{lemma}
\begin{proof}
Call a target active when $p_x>0$.
Starting from a unit-slack target with $p_x>0$, iterate
Lemma~\ref{lem:parent}. The strictly increasing $q$ values force
termination at a higher-slack target.
Such a starting target exists unless the required higher-slack targets
already exist, because minimum-degree sources have far targets in $K$.

Suppose that the only higher-slack target is $c$. Every active
unit-slack target $x$ then satisfies
$P_x\cup\{x\}\subseteq P_c$.
A minimum-degree source has at least two far targets in $K$,
so one differs from $c$; it follows that the source lies in $P_c$.
Every nonminimum-degree vertex $v\ne c$ is unit-slack and has $e_v=1$,
hence $|U(v)|=g(v)\ge1$.
It cannot be far from a zero-slack target, whose far sources have
minimum degree. Its far target is therefore either $c$ or an active
unit-slack target; in both cases $v\in P_c$.
Thus $P_c=V(D)\setminus\{c\}$, which makes the indegree of $c$ zero,
contradicting strong connectivity.
The case of no higher-slack target is already excluded by the
strictly increasing chains.
\end{proof}

Every unit-slack vertex has an incoming arc from $K$.
For $p_x>0$ this is Lemma~\ref{lem:parent}. If $p_x=0$, then $q_x=1$.
When $e_x=1,\mu(x)=0$, a far target in $K$ exists and is adjacent to $x$,
so points to $x$. When $e_x=0,\mu(x)=1$, at least two far targets lie
in $K$, at most one is missing from $x$, and an adjacent one again
points to $x$. Distinct unit-slack vertices give arcs with distinct
heads.

\subsection{The missing-pair balance inside $K$}

\begin{theorem}\label{thm:wholecore}
In the setting of this section,
\[
{z\ge h+7.}
\]
Moreover $k\ge5$ and
$k\le\lfloor(n-h-2)/2\rfloor$.
\end{theorem}
\begin{proof}
Let $\ell$ be the number of unit-slack targets and $H=k-\ell$ the
number of higher-slack targets. Let $A_K$ count arcs inside $K$,
let $m$ count minimum-degree vertices in $K$, let $t$ count
gap-two vertices with excess at most one, and let $M_{\rm out}$
count missing pairs with both endpoints outside $K$.
Thus $h=r_2+t$.

By Lemma~\ref{lem:zeroslack}, outside $K$ there are $z$ whole
minimum-degree vertices and $L=n-k-z$ minimum-degree vertices of
missing degree one. Their missing-degree sum is
$L=M(K,V\setminus K)+2M_{\rm out}$.
As $r_0=n-k+m$ and $M=r_0-r_2$, the number of missing pairs
inside $K$ equals
\[
M_K=M-L+M_{\rm out}
=m-r_2+z+M_{\rm out}
=\Sigma+m+t+M_{\rm out}.
\]
Therefore
\begin{equation}\label{eq:corebalance}
{\binom{k}{2}=\Sigma+m+t+A_K+M_{\rm out}.}
\end{equation}
The preceding incoming-arc argument gives $A_K\ge\ell$, while
$\Sigma+\ell\ge2k$.
Hence $\binom{k}{2}\ge2k$, and $k\ge5$ since $k>0$.
Lemma~\ref{lem:twoheavy} gives $H\ge2$, so
$\Sigma\ge k+H\ge7$. Equation~\eqref{eq:slacksum} proves $z\ge h+7$.
Finally, $z\le n-k$ and $\Sigma\ge k+2$ give
$n\ge 2k+h+2$, proving the upper bound on $k$.
\end{proof}

In particular, at least seven minimum-degree vertices are adjacent to
every other vertex. This does not exclude the order $2\delta+3$ itself;
it supplies additional structure that any such minimal counterexample
must possess.
